% !TeX root = ../main.tex


\section*{Заключение}

Стандартная оценка контраста с канонической HRF при нелинейной гемодинамике смещена, и смещение зависит от амплитуд, плотности и последовательности событий. Собственное линейное ядро генератора устраняет лишь ошибку формы HRF и этого смещения не снимает, поскольку нелинейный вклад остается. Несовпадение канонической формы может как усиливать, так и частично компенсировать нелинейный вклад, поэтому близость форм сама по себе точности не гарантирует.

Ранее нелинейность BOLD-отклика подтверждалась качественно, по соответствию модельных кривых наблюдаемому сигналу. В настоящей работе ошибки активационных оценок вычислены против точного эталона и разложены на вклад формы HRF, нелинейной динамики и шума, а границы применимости линейного анализа и фиксированных поправок определены количественно, по заранее заданным порогам. Приближение второго порядка верно воспроизводит знак ошибки контраста и потому пригодно для диагностики ее направления.

Для практики это означает, что при плотных последовательностях событий контраст нельзя рассматривать как прямую меру разности амплитуд воздействия, а различия в последовательностях предшествующих событий способны давать ненулевую оценку при равных амплитудах. Численные границы относятся к исследованным моделям и протоколам, их перенос на данные человека требует отдельного обоснования.
